\documentclass[11pt]{article}
\usepackage[T1]{fontenc}
\usepackage{lmodern,amsmath,amssymb,amsthm,mathtools}
\usepackage[margin=1in]{geometry}
\usepackage{microtype,enumitem,xcolor}
\usepackage[numbers,sort&compress]{natbib}
\usepackage{hyperref}
\hypersetup{colorlinks=true,linkcolor=blue!45!black,citecolor=blue!45!black,
urlcolor=blue!45!black,pdftitle={Zero-bias transport and extremal Bernoulli arrays},pdfauthor={Logan Bell}}
\setlength{\emergencystretch}{2em}
\numberwithin{equation}{section}
\newtheorem{theorem}{Theorem}[section]
\newtheorem{proposition}[theorem]{Proposition}
\newtheorem{lemma}[theorem]{Lemma}
\theoremstyle{remark}
\newtheorem{remark}[theorem]{Remark}
\DeclareMathOperator{\Var}{Var}
\DeclareMathOperator{\Bern}{Bern}
\DeclareMathOperator{\Bin}{Bin}
\DeclareMathOperator{\supp}{supp}
\newcommand{\E}{\mathbb E}
\newcommand{\Pp}{\mathbb P}
\newcommand{\R}{\mathbb R}
\newcommand{\Z}{\mathbb Z}
\newcommand{\ce}{C_{\mathrm E}}
\newcommand{\phz}{\phi(0)}
\newcommand{\Law}{\mathcal L}
\title{Zero-bias transport and extremal Bernoulli arrays}
\author{Logan Bell}
\date{6 October 2026}
\begin{document}
\maketitle
\begin{abstract}
For a finite sum of independent centered real variables, the Wasserstein
distance to its zero-biased law is at most the sum of the third absolute
moments divided by twice the total variance. We prove that equality holds
exactly when all nonconstant summands have one identical two-point law.
The proof treats arbitrary summand distributions: a same-variance Choquet
representation reduces the single-law equality problem to three-point
laws, and an exact conditional-variance identity resolves the remaining
Bernoulli mixtures. A sharp comparison between lattice span, skewness,
and zero-bias transport then identifies equality in the Berry--Esseen
limit obtained by independently repeating a fixed finite block. The
limiting ratio is at most Esseen's constant, with equality exactly for
identically distributed Esseen two-point summands and any zero summands.
\end{abstract}

\setcounter{tocdepth}{2}
\tableofcontents

\section{The problem and the main results}\label{sec:results}
The classical extremal example for the Berry--Esseen inequality is a sum
of identically distributed two-point variables. For general independent
summands, two separate questions arise: what forces two-point support,
and what forces the resulting laws to agree? This paper answers both
questions for equality in the optimal zero-bias transport bound. It then
uses that equality classification in a fixed-block Berry--Esseen limit.

For probability laws with finite first moments, write
\[
 W_1(\mu,\nu)=\inf\{\E|U-V|:\Law(U)=\mu,\ \Law(V)=\nu\}.
\]
We also write $W_1(U,V)$ for the distance between their laws. If $X$ is
centered and has variance $v>0$, its zero-biased law is characterized by
\begin{equation}\label{eq:zero-bias}
             \E[Xf(X)]=v\E f'(X^*).
\end{equation}
This identity holds for every absolutely continuous $f$ for which the
expectations exist.
Goldstein and Reinert introduced this transformation and its
independent-sum replacement construction
\cite[Definition~1.1 and Lemma~2.1]{GoldsteinReinert1997}.
Goldstein and Tyurin independently proved the sharp bound
\begin{equation}\label{eq:single-bound}
 W_1(X,X^*)\le\frac{\E|X|^3}{2v}
\end{equation}
for every centered law with finite third absolute moment
\cite[Lemma~1.1]{Goldstein2010}\cite[Theorem~3]{Tyurin2009}.
Every nonconstant two-point law attains equality. Our first result
characterizes all equality cases, both for one law and for an independent
array.

Let $X_1,\ldots,X_n$ be independent centered real variables and put
\begin{equation}\label{eq:array-data}
 S=\sum_{i=1}^nX_i,\qquad
 V=\sum_{i=1}^n\E X_i^2>0,\qquad
 R=\sum_{i=1}^n\E|X_i|^3<\infty.
\end{equation}
The denominator $R$ records the original summands. It is generally
different from $\E|S|^3$.

\begin{theorem}[Equality in zero-bias transport]\label{thm:array-equality}
For the array in \eqref{eq:array-data},
\begin{equation}\label{eq:additive-bound}
                      2V W_1(S,S^*)\le R.
\end{equation}
Equality holds if and only if all nonconstant coordinates have one
identical centered two-point law. Any remaining coordinates are
identically zero. In particular, equality in \eqref{eq:single-bound}
holds if and only if $X$ has a two-point law.
\end{theorem}

The distinction between $R$ and $\E|S|^3$ is necessary even for the
equality statement. A sum of two or more nonconstant iid Bernoulli
variables has more than two support points. Its single-law bound is
therefore strict, although its bound with the original third-moment
sum is an equality.

For changing arrays with Lyapunov ratio tending to zero and
Berry--Esseen ratio tending to Esseen's constant, a companion paper
\cite{BellSmallStructure2026} proves that \eqref{eq:additive-bound}
approaches equality. Full finite-horizon global optimality in that
regime gives the exact equality class. The fixed-block theorem below
has its own asymptotic hypotheses.

The second result concerns repetitions of one fixed independent block.
Let $S_1,S_2,\ldots$ be independent copies of the sum $S$ in
\eqref{eq:array-data}. For each integer $m\ge1$, define
\begin{equation}\label{eq:block-ratio}
 C_m=\frac{\sqrt m\,V^{3/2}}{R}
 \sup_x\left|\Pp\left\{\frac{S_1+\cdots+S_m}{\sqrt{mV}}\le x\right\}
                   -\Phi(x)\right|.
\end{equation}
This is the Berry--Esseen ratio for the $mn$ original coordinates,
whose total variance is $mV$ and whose third-moment sum is $mR$.
Set
\[
 \ce=\frac{3+\sqrt{10}}{6\sqrt{2\pi}},\qquad
 p_{\mathrm E}=\frac{4-\sqrt{10}}2.
\]

\begin{theorem}[A fixed block repeated independently]\label{thm:block}
The limit of \eqref{eq:block-ratio} exists and satisfies
\begin{equation}\label{eq:block-bound}
 \lim_{m\to\infty}C_m
 \le \ce\frac{2V W_1(S,S^*)}{R}\le\ce.
\end{equation}
Equality with $\ce$ holds if and only if, for some $h>0$, every
nonconstant coordinate has the law of $h(B-p_{\mathrm E})$, where
$B\sim\Bern(p_{\mathrm E})$, or every such coordinate has its reflection.
The other coordinates are zero.
\end{theorem}

Esseen's fixed-law asymptotic formula and sharp moment inequality already
give the constant $\ce$ for iid sums \cite[pp.~160--163]{Esseen1956}.
Theorem~\ref{thm:block} retains the separate third moments of the original
block and identifies equality through transport. The proof depends on a
lattice inequality in which entire binomial laws, rather than only
two-point laws, attain equality. The original denominator then determines
which binomial decompositions attain the Berry--Esseen constant.

The block is fixed throughout the limit. The theorem gives no uniform
remainder estimate for changing laws, and it does not prove the sharp
Berry--Esseen bound at a finite sample size. Its proof uses the classical
asymptotic formula.

\section{Zero bias and the independent-sum bound}\label{sec:basic}
We use two formulas from Goldstein and Reinert
\cite[Lemma~2.1(2), (4), and (5)]{GoldsteinReinert1997}. The zero-biased
law of a centered variable $X$ of variance $v>0$ has density
\begin{equation}\label{eq:density}
             d_X(u)=\frac1v\E[X\mathbf1_{\{X>u\}}].
\end{equation}
For $u<0$, centering rewrites the numerator as
$-\E[X\mathbf1_{\{X\le u\}}]$, so the density is nonnegative on both
half-lines. The defining identity gives
\begin{equation}\label{eq:bias-moments}
 \E X^*=\frac{\E X^3}{2v},\qquad
 \E|X^*|=\frac{\E|X|^3}{2v}.
\end{equation}
The first follows by taking $f(x)=x^2/2$; the second follows by taking
$f(x)=x|x|/2$.

For laws with finite first moments, the Wasserstein distance is the
integral of the absolute CDF difference:
\begin{equation}\label{eq:cdf-distance}
       W_1(\mu,\nu)=\int_\R|F_\mu(x)-F_\nu(x)|\,dx.
\end{equation}
For example, Goldstein states the optimal quantile coupling in
\cite[equations~(15)--(16)]{Goldstein2010}. This formula makes both mixture
convexity and convolution contraction immediate by the triangle
inequality and Tonelli's theorem.

Discard the zero coordinates of \eqref{eq:array-data} and write
$v_i=\E X_i^2$, $\rho_i=\E|X_i|^3$, and $T_i=S-X_i$. Choose an index
$I$, independently of all variables, with
$\Pp(I=i)=v_i/V$, and replace $X_I$ by an independent variable with its
zero-biased law. Then
\begin{equation}\label{eq:replacement}
                  S^*\ \stackrel{d}=\ T_I+X_I^*.
\end{equation}
Indeed, applying \eqref{eq:zero-bias} to each summand, conditional on
the others, gives
$\E[Sf(S)]=\sum_i v_i\E f'(T_i+X_i^*)$.
Consequently
\begin{equation}\label{eq:bound-chain}
 \begin{split}
 W_1(S,S^*)
 &\le\sum_i\frac{v_i}{V}W_1(S,T_i+X_i^*)\\
 &\le\sum_i\frac{v_i}{V}W_1(X_i,X_i^*)
 \le\frac R{2V}.
 \end{split}
\end{equation}
The last step is \eqref{eq:single-bound}. Tyurin's Corollary~1 gives
the same independent-sum bound \cite{Tyurin2009}.

There are three possible sources of strictness in
\eqref{eq:bound-chain}: a strict single-coordinate bound, cancellation
under convolution, and cancellation between replacement components.
The next two sections determine their equality cases.

\section{Equality for an arbitrary single law}\label{sec:single-equality}
For a centered law $\mu$ of variance $v>0$ and third absolute moment
$\rho$, put
\[
                 B(\mu)=\frac{2v W_1(\mu,\mu^*)}{\rho}.
\]
This quantity is invariant under nonzero scaling and reflection. We first
settle laws with at most three atoms, then use an exact mixture at fixed
variance to treat every law.

\subsection{Two-point laws and strict three-point bounds}
A centered two-point law can be written as $h(B-p)$, with
$h>0$, $0<p<1$, $q=1-p$, and $B\sim\Bern(p)$. Its variance and third
absolute moment are
\[
        v=h^2pq,\qquad \rho=h^3pq(p^2+q^2).
\]
Its zero-biased law is uniform on $[-hp,hq]$, and direct integration
gives
\begin{equation}\label{eq:binary-transport}
                 W_1(\mu,\mu^*)=\frac h2(p^2+q^2).
\end{equation}
These are Goldstein's exact two-point formulas
\cite[Lemma~1.2]{Goldstein2010}; they give $B(\mu)=1$.

\begin{lemma}\label{lem:strict-ternary}
Every centered law with three distinct support points of positive mass
has $B(\mu)<1$.
\end{lemma}
\begin{proof}
First suppose none of its atoms is zero. Reflect the law, if necessary,
so that its support is $x<y<0<z$. Let $\mu_1$ and $\mu_0$ be the
centered two-point laws on $\{x,z\}$ and $\{y,z\}$. There is a unique
$\alpha\in(0,1)$ such that
\[
                 \mu=\alpha\mu_1+(1-\alpha)\mu_0.
\]
The variances are $-xz$ and $-yz$. Thus
\[
 \mu^*=\beta\mu_1^*+(1-\beta)\mu_0^*,\qquad
 \beta=\frac{\alpha x}{\alpha x+(1-\alpha)y}>\alpha.
\]
This variance reweighting is Goldstein's mixture identity
\cite[Theorem~2.1]{Goldstein2010}.

Write $D_j=W_1(\mu_j,\mu_j^*)$ and
$D_{10}=W_1(\mu_0,\mu_1^*)$. A coupling formed with weights
$\alpha$, $1-\beta$, and $\beta-\alpha$ gives
\begin{equation}\label{eq:ternary-coupling}
 W_1(\mu,\mu^*)
 \le\alpha D_1+(1-\beta)D_0+(\beta-\alpha)D_{10}.
\end{equation}
In contrast, the two-point equalities give
\[
                   \frac\rho{2v}=\beta D_1+(1-\beta)D_0.
\]
The cross-transport comparison in Goldstein's Lemma~3.1 is strict for
$x<y<0<z$. More explicitly, integrating the two affine CDF pieces,
as in \cite[equation~(28) and the following case]{Goldstein2010}, gives
\[
 D_1-D_{10}=
 \begin{cases}
 \displaystyle
 \frac{-y(y-x)\{y^2+2z^2-y(x+2z)\}}{(z-x)(z-y)^2},
       &y(x+z)\le y^2+z^2,\\[2ex]
 \displaystyle\frac{z^2}{z-x},&y(x+z)>y^2+z^2.
 \end{cases}
\]
In the first case the expression in braces is at least $z(z-y)>0$.
The second expression is positive as well. Subtracting
\eqref{eq:ternary-coupling} therefore gives a strictly positive deficit
$\rho/(2v)-W_1(\mu,\mu^*)$.

It remains to treat an atom at zero. Write the support as
$\{-ph,0,qh\}$, where $p+q=1$ and, after reflection, $0<p\le q$.
If the middle mass is $1-\tau\in(0,1)$, the endpoint masses are
$\tau q$ and $\tau p$. The zero-biased law is uniform on the full
interval. Integration on its two sides gives, with $r=p^2+q^2$,
\begin{equation}\label{eq:zero-middle}
 B(\mu)=
 \begin{cases}
 (q-p+2\tau^2p^2)/r,&\tau q\ge p,\\
 1+2\tau^2-4\tau pq/r,&\tau q<p.
 \end{cases}
\end{equation}
The deficits are respectively
$2p^2(1-\tau^2)/r$ and $2\tau(2pq/r-\tau)$. Both are positive:
in the second case, $2pq/r\ge p/q>\tau$. This direct calculation
provides strictness at the boundary which Goldstein's proof of
Theorem~3.1 treats by continuity.
\end{proof}

\subsection{An exact mixture at fixed variance}\label{sec:choquet}
Normalize the variance of $\mu$ to one. Consider the convex set
\[
 \mathcal K=\left\{\nu:\nu\text{ is a probability law on }\R,
         \ \int x\,\nu(dx)=0,\ \int x^2\,\nu(dx)\le1\right\}.
\]
The inequality on the second moment makes this set weakly compact.
Indeed, it gives tightness and uniform integrability of first moments,
since
\[
           \int_{|x|>A}|x|\,\nu(dx)\le\frac1A.
\]
Centering is therefore preserved under weak limits, and lower
semicontinuity preserves the second-moment bound. The weak topology on
$\mathcal K$ is metrizable. Embed finite signed measures in the locally
convex Hausdorff space whose topology is induced by integration against
$C_b(\R)$. Choquet's representation theorem applies to $\mathcal K$
\cite[pp.~333--337, especially Corollary~2 and Lemma~5]{Choquet1960}.
It gives a probability measure $\pi$ carried by the extreme points,
with barycenter $\mu$.

Every extreme point of $\mathcal K$ has at most three atoms. To see
this, suppose $\nu$ charges four disjoint Borel sets $A_1,\ldots,A_4$.
Choose a nonzero vector $(c_j)_{j=1}^4$ in the nullspace of their three
feature rows
\[
        \left(\int_{A_j}1\,d\nu\right)_j,\quad
        \left(\int_{A_j}x\,d\nu\right)_j,\quad
        \left(\int_{A_j}x^2\,d\nu\right)_j.
\]
The bounded function $g=\sum_jc_j\mathbf1_{A_j}$ is nonzero
$\nu$-almost surely on a set of positive mass. For sufficiently small
$\varepsilon>0$, the distinct measures $(1\pm\varepsilon g)\nu$
belong to $\mathcal K$ and have average $\nu$, contradicting
extremality. A law not supported on at most three points admits such
four sets.

The barycenter identity first holds for bounded continuous functions.
It identifies $\mu$ with the mixture measure $\int\nu\,\pi(d\nu)$,
and hence extends to nonnegative Borel functions. In particular,
\begin{equation}\label{eq:exact-variance-mixture}
 1=\int\left(\int x^2\,\nu(dx)\right)\pi(d\nu),\qquad
 \rho(\mu)=\int\rho(\nu)\,\pi(d\nu).
\end{equation}
Every inner second moment is at most one. Therefore it equals one
for $\pi$-almost every component. The decomposition is consequently
into centered laws of the \emph{same} variance, with no degenerate
component.

For these components the zero-bias transformation uses the same
mixture weights. One can check this without any choice of coupled
variables: for a Borel set $A$, the measurable positive kernel
\begin{equation}\label{eq:bias-kernel}
             \nu^*(A)=\int x\left(\int_0^x\mathbf1_A(u)\,du\right)\nu(dx)
\end{equation}
has total mass one and satisfies \eqref{eq:zero-bias}. For negative
$x$, the inner integral is oriented, so its product with $x$ remains
nonnegative. Tonelli's theorem gives
\begin{equation}\label{eq:same-weight}
                        \mu^*=\int\nu^*\,\pi(d\nu).
\end{equation}
All required first moments are integrable by
\eqref{eq:bias-moments} and \eqref{eq:exact-variance-mixture}. Thus
\begin{equation}\label{eq:mixture-bound}
 W_1(\mu,\mu^*)
 \le\int W_1(\nu,\nu^*)\,\pi(d\nu)
 \le\frac12\int\rho(\nu)\,\pi(d\nu)=\frac{\rho(\mu)}2.
\end{equation}
Goldstein's Theorem~2.2 gives this convexity at fixed variance; it also
explains why unequal mixture variances cannot be used here
\cite[Remark~2.2]{Goldstein2010}. The exact barycenter above permits
its use in the equality problem for an arbitrary law.

If $B(\mu)=1$, the nonnegative component deficits in
\eqref{eq:mixture-bound} have integral zero. Lemma~\ref{lem:strict-ternary}
then shows that almost every component is a two-point law.

\subsection{Distinct binary components cancel}
A centered unit-variance two-point law is uniquely parametrized by
$a>0$:
\[
 \nu_a=\frac1{1+a^2}\delta_{-a}
              +\frac{a^2}{1+a^2}\delta_{1/a}.
\]
Its signed third moment is $c(a)=1/a-a$, a strictly decreasing
continuous function with a continuous inverse. This parametrization
is Borel, since moments are measurable limits of bounded truncations.
Push the representing measure $\pi$ to the parameter $a$.

The signed CDF difference from the uniform zero-biased law is
\begin{equation}\label{eq:binary-cdf}
 f_a(x)=F_{\nu_a}(x)-F_{\nu_a^*}(x)
 =\frac1{1+a^2}-\frac{x+a}{a+1/a}\quad(-a<x<1/a),
\end{equation}
and is zero outside that interval, apart from irrelevant endpoint
values. Its zero is $c(a)$. If $a<b$, then $f_a>0$ and $f_b<0$ on
\begin{equation}\label{eq:opposite-interval}
          (\max\{-a,c(b)\},\ \min\{1/b,c(a)\}).
\end{equation}
This interval has positive length: each of $-a$ and $c(b)$ is
strictly smaller than each of $1/b$ and $c(a)$.

If the parameter distribution is not a point mass, choose two
distinct points $a<b$ in its support. On a compact subinterval of
\eqref{eq:opposite-interval}, sufficiently small parameter
neighborhoods retain the two opposite signs, bounded away from zero.
Both neighborhoods have positive probability. For every $x$ in that
subinterval,
\[
 \E|f_a(x)|-|\E f_a(x)|
       =2\min\{\E(f_a(x))_+,\E(f_a(x))_-\}>0.
\]
Here $u_+=\max(u,0)$ and $u_-=\max(-u,0)$. Integrating gives a
strict triangle inequality in the first line of
\eqref{eq:mixture-bound}. Fubini is justified by
$\int\E|f_a(x)|\,dx=\int\rho(\nu_a)\,\pi(da)/2<\infty$.
Therefore equality forces a single binary component. Scaling back
proves the single-law assertion in Theorem~\ref{thm:array-equality}.

\section{Equality for independent arrays}\label{sec:array}
If any nonconstant coordinate is not two-point, the last inequality
in \eqref{eq:bound-chain} is strict. We may therefore write every
nonconstant coordinate as
\[
                  X_i=h_i(B_i-p_i),\qquad h_i>0,\quad 0<p_i<1,
\]
with independent $B_i\sim\Bern(p_i)$. Two remaining calculations
determine the gaps and probabilities separately.

\subsection{Unequal support diameters give strict convolution loss}
Suppose a coordinate has the largest support diameter $h$ and another
coordinate has diameter $0<\delta<h$. In raw coordinates the signed
CDF difference between the first coordinate and its uniform zero-biased
law is
\[
                f(x)=q-x/h\quad(0<x<h),\qquad f(x)=0\text{ otherwise},
\]
where $q$ is its lower-atom probability. The deleted sum has two atoms
of positive probability separated by $\delta$: fix all its Bernoulli
coordinates at their lower atoms, and then change only the coordinate
of diameter $\delta$.

The two corresponding translates of $f$ have opposite signs on a
translate of
\[
                    (\max\{qh,\delta\},\ \min\{qh+\delta,h\}).
\]
Its length is positive because $0<q<1$ and $0<\delta<h$.
Consequently $|\E f(x-T)|<\E|f(x-T)|$ on a set of positive measure.
After integration, convolution contraction is strict for the
largest-gap coordinate in \eqref{eq:bound-chain}. Equality in
\eqref{eq:additive-bound} therefore forces one common diameter $h$.

\subsection{An exact deficit for a common diameter}\label{sec:posterior}
Assume there are $n$ nonconstant coordinates, all of diameter $h$.
Put
\[
 K=\sum_i B_i,\qquad \mu=\sum_i p_i,\qquad
 w_i=\frac{h^2p_i(1-p_i)}V,\qquad
 c_{i,k}=\Pp\left\{\sum_{j\ne i}B_j=k\right\}.
\]
For $0\le k\le n-1$, define
\begin{equation}\label{eq:posterior-weights}
 b_k=\sum_i w_i c_{i,k},\qquad
 \pi_{i,k}=\frac{w_i c_{i,k}}{b_k},\qquad
 Q_k=\sum_i\pi_{i,k}(1-p_i).
\end{equation}
Every $b_k$ is positive, and $\sum_kb_k=1$.

\begin{proposition}[Conditional-variance identity]\label{prop:posterior}
For a common-diameter Bernoulli array,
\begin{equation}\label{eq:posterior-deficit}
 \frac R{2V}-W_1(S,S^*)
       =h\sum_{k=0}^{n-1}b_k\Var_{\pi_k}(p_i)
       =\frac h2\sum_{k=0}^{n-1}
          \frac{\sum_{i,j}w_iw_jc_{i,k}c_{j,k}(p_i-p_j)^2}{b_k}.
\end{equation}
\end{proposition}
\begin{proof}
On the sum-lattice cell $[h(k-\mu),h(k+1-\mu)]$, let $s\in(0,1)$
be the relative position. Condition on the deleted count. The
$i$th replacement in \eqref{eq:replacement} gives exactly
\[
       F_S-F_{T_i+X_i^*}=c_{i,k}(1-p_i-s).
\]
After averaging with the variance weights, this becomes
$F_S-F_{S^*}=b_k(Q_k-s)$. The integral of its absolute value on the
cell is
\[
                 \frac{h b_k}2\{Q_k^2+(1-Q_k)^2\}.
\]
On the other hand, \eqref{eq:binary-transport} gives
\[
 \frac R{2V}=\frac h2\sum_iw_i\{p_i^2+(1-p_i)^2\}
            =\frac h2\sum_kb_k\E_{\pi_k}\{p_i^2+(1-p_i)^2\}.
\]
Subtracting yields the first identity, since the difference between
the two quadratic expressions is $2\Var_{\pi_k}(p_i)$. The usual
pairwise formula for variance gives the second.
\end{proof}

The weights $\pi_{i,k}$ have a direct probabilistic meaning: conditional
on the chosen zero-bias replacement interval having cell index $k$,
they are the probabilities of the replaced coordinate. At $k=0$ every
coordinate has positive conditional weight. Thus
\eqref{eq:posterior-deficit} vanishes exactly when all $p_i$ agree.
Conversely, for an iid binary sum all the cell means are $Q_k=1-p$,
and
\begin{equation}\label{eq:iid-transport}
                      W_1(S,S^*)=\frac h2\{p^2+(1-p)^2\}.
\end{equation}
Together with the strict unequal-diameter argument, this proves
Theorem~\ref{thm:array-equality}. Zero coordinates have zero variance
weight and do not affect any of the identities.

\section{Lattice span, skewness, and transport}\label{sec:lattice}
We now compare the two terms in Esseen's fixed-law asymptotic formula
with the transport distance. Let $X$ be a centered lattice variable
of variance $v>0$ and finite third absolute moment. Suppose its law
is supported on $a+h\Z$, with $h>0$, and write $\kappa=\E X^3$.
The spacing $h$ need not be maximal for the inequality below.

\begin{theorem}[A lattice transport inequality]\label{thm:span}
For such a law,
\begin{equation}\label{eq:span-transport}
              |\kappa|+3hv\le(3+\sqrt{10})\,2v W_1(X,X^*).
\end{equation}
Equality holds exactly for the laws of
$h(N-np)$, where $N\sim\Bin(n,p)$, $n\ge1$, and
$p\in\{p_{\mathrm E},1-p_{\mathrm E}\}$.
\end{theorem}

\subsection{Integration on the lattice cells}
By \eqref{eq:density}, the density of $X^*$ is constant on every open
cell $(a+kh,a+(k+1)h)$; call it $d_k$. If $d_k=0$, the cell lies
outside the support's convex hull and its CDF difference is zero.
For the other cells put
\[
 \omega_k=hd_k,\qquad
 r_k=\frac{F_X(a+kh)-F_{X^*}(a+kh)}{hd_k},\qquad
 \ell(r)=\int_0^1|r-s|\,ds.
\]
Then $\sum_k\omega_k=1$. The CDF difference within the cell is
$hd_k(r_k-s)$, so
\begin{equation}\label{eq:cell-identities}
 \frac{W_1(X,X^*)}h=\sum_k\omega_k\ell(r_k),\qquad
 \sum_k\omega_k(r_k-1/2)=\frac\kappa{2vh}.
\end{equation}
For the second identity, integrate the signed CDF difference and use
$\E X^*-\E X=\kappa/(2v)$. Absolute convergence follows from the
finite first moments in \eqref{eq:bias-moments}.

The function $\ell$ is convex, with
\[
 \ell(r)=
 \begin{cases}
 r^2-r+1/2,&0\le r\le1,\\
 |r-1/2|,&r\notin[0,1].
 \end{cases}
\]
Let $t=|\kappa|/(vh)$. Jensen's inequality in
\eqref{eq:cell-identities} gives
\begin{equation}\label{eq:jensen-bound}
 2vW_1(X,X^*)\ge
 \begin{cases}
 vh(1+t^2)/2,&0\le t\le1,\\
 vht,&t\ge1.
 \end{cases}
\end{equation}
In the first range, $2(3+t)/(1+t^2)$ has its unique maximum
$3+\sqrt{10}$ at $t=\sqrt{10}-3$. In the second range,
$(3+t)/t\le4<3+\sqrt{10}$. This proves
\eqref{eq:span-transport}.

Combining \eqref{eq:span-transport} with
\eqref{eq:single-bound} recovers Esseen's lattice moment inequality
\cite[Lemma~1, pp.~162--163]{Esseen1956}:
\[
                |\kappa|+3hv\le(3+\sqrt{10})\E|X|^3.
\]
Esseen's original lemma allows the more general condition that consecutive
support points are at least $h$ apart. We use the lattice case because
the asymptotic correction is determined by lattice span.

\subsection{Equality determines a binomial law}
Equality in \eqref{eq:span-transport} requires
$|\kappa|/(vh)=\sqrt{10}-3<1$. The mean argument of $\ell$ in
Jensen's inequality therefore lies strictly between zero and one.
At that argument its supporting line touches $\ell$ only once.
All positive-density cells consequently have the same ratio $r$, with
\begin{equation}\label{eq:r-optimal}
                 0<r<1,\qquad |2r-1|=\sqrt{10}-3.
\end{equation}

Let $x$ be a grid point having mass $p_x$, and let $d_-$ and $d_+$
be the densities on its left and right cells. The CDF jump and the
density jump give
\[
 p_x=h\{r d_++(1-r)d_-\},\qquad
                   v(d_--d_+)=x p_x.
\]
These equations remain valid at the support endpoints, where one
density is zero. Solving them gives
\begin{equation}\label{eq:binomial-density}
 hd_-=p_x(1+rhx/v),\qquad
 hd_+=p_x(1-(1-r)hx/v).
\end{equation}
A missing interior grid atom is impossible: its positive unchanged
density would make the cell ratio decrease by one across that point,
contrary to its constancy. Positivity in
\eqref{eq:binomial-density} confines every atom to
$[-v/(rh),v/((1-r)h)]$, so the support is finite. At its lower and
upper endpoints,
\[
                   x_-=-\frac v{rh},\qquad
                   x_+=\frac v{(1-r)h}.
\]
Their separation is $nh$ for an integer $n\ge1$. Thus, with $p=1-r$,
\[
          v=nr(1-r)h^2,\qquad x_j=h(j-np),\quad 0\le j\le n.
\]
The density on the cell between $x_j$ and $x_{j+1}$ is the same
when calculated from either endpoint. Substituting these locations
in \eqref{eq:binomial-density} gives
\[
             \frac{p_{x_{j+1}}}{p_{x_j}}
                    =\frac{(n-j)p}{(j+1)(1-p)}.
\]
Normalization identifies the law as centered $h\Bin(n,p)$, and
\eqref{eq:r-optimal} gives $p=p_{\mathrm E}$ or $1-p_{\mathrm E}$.
Conversely, \eqref{eq:iid-transport} and the binomial third moment
verify equality for each of these laws. This completes the proof of
Theorem~\ref{thm:span}.

For $n\ge2$ these binomial laws are not two-point laws. They attain
the span/transport inequality, but their single-law transport bound
is strict. Resolving them into their independent binary constituents
changes the third-moment denominator and restores equality in
\eqref{eq:additive-bound}.

\section{The Berry--Esseen limit for a fixed block}\label{sec:block}
Return to \eqref{eq:array-data}. The law of $S$ has finite third
absolute moment, since the block is finite. Its variance is $V$ and
its signed third moment is
\[
                          \kappa=\E S^3=\sum_i\E X_i^3.
\]
Let $h$ be its maximal lattice span, or set $h=0$ if its law is
nonlattice. Esseen's fixed-law asymptotic formula
\cite[pp.~160--162, equations~(6)--(7)]{Esseen1956} gives
\begin{equation}\label{eq:esseen-limit}
 \lim_{m\to\infty}\sqrt m\,
 d_K\left(\frac{S_1+\cdots+S_m}{\sqrt{mV}},Z\right)
                  =\frac{\phz}{6V^{3/2}}(|\kappa|+3hV),
\end{equation}
where $Z$ is standard normal and $d_K$ is Kolmogorov distance.
The source uses one fixed law with positive variance and finite third
absolute moment. Multiplying by $V^{3/2}/R$ yields
\begin{equation}\label{eq:original-price-limit}
                    \lim_{m\to\infty}C_m
                      =\frac{\phz}{6R}(|\kappa|+3hV).
\end{equation}

If $h>0$, Theorem~\ref{thm:span} applies to the sum law and gives
the first inequality in \eqref{eq:block-bound}. Notice that
\eqref{eq:span-transport} has no third absolute moment in its
denominator, so it can be divided by the original $R$ without
identifying $R$ with $\E|S|^3$. The second inequality is
Theorem~\ref{thm:array-equality}.

If $h=0$, the transport distance dominates the difference of means,
so \eqref{eq:bias-moments} gives
$|\kappa|\le2V W_1(S,S^*)$. Hence the stronger comparison
\[
             \lim_{m\to\infty}C_m
             \le\frac{\phz}{6}\frac{2V W_1(S,S^*)}{R}
             \le\frac{\phz}{6}<\ce
\]
holds in the nonlattice case.

Finally suppose the limit equals $\ce$. The sum is lattice, and
both inequalities in \eqref{eq:block-bound} are equalities.
Theorem~\ref{thm:array-equality} forces all nonconstant coordinates
to have one law $g(B-p)$, with $g>0$. If there are $k$ such
coordinates, the sum is centered $g\Bin(k,p)$; its maximal lattice
span is $g$, and
\[
                        \frac\kappa{Vg}=1-2p.
\]
Equality in Theorem~\ref{thm:span} forces
$|1-2p|=\sqrt{10}-3$, or $p\in\{p_{\mathrm E},1-p_{\mathrm E}\}$.
Conversely these arrays attain both equalities and have the stated
limit by \eqref{eq:esseen-limit}. This proves
Theorem~\ref{thm:block}.

\bibliographystyle{plainnat}
\bibliography{refs}
\end{document}
